\section{Chronological Development}
\label{sec:chronology}

The four clusters did not mature in lockstep. Strong consistency arrived first, borrowing directly from
data-centre consensus literature; eventual consistency followed as mobile and IoT deployments made
partitions routine rather than exceptional; tunable models emerged once enough operational data existed
to show neither extreme fit every workload; and placement-aware co-design is the most recent and least
settled thread. Table~\ref{tab:timeline} lists one representative milestone per year where the corpus
contains a clear inflection point.

\begin{table}[h!]
\centering
\small
\begin{tabularx}{\textwidth}{@{}p{1.6cm}Xp{2.6cm}@{}}
\toprule
\rowcolor{nx@teal}
\textcolor{white}{\textbf{Year}} & \textcolor{white}{\textbf{Milestone}} & \textcolor{white}{\textbf{Cluster}} \\
\midrule
2015 & \citet{chen2015nexa} formalise lease-based invalidation bounds under wireless jitter (Nexa). & A. Strong \\
\rowcolor{nx@mist}
2017 & \citet{patel2017synapse} adapt CRDT merges to asymmetric edge bandwidth (Synapse). & B. Eventual \\
2018 & First large-scale field study quantifies staleness tolerance across five application classes, motivating tunable bounds. & C. Tunable \\
\rowcolor{nx@mist}
2019 & \citet{osei2019vertex} expose staleness as a per-tenant scheduling parameter (Vertex). & C. Tunable \\
2021 & \citet{alvarez2021nimbus} show placement-only optimisers under-provision write-heavy replica sets (Nimbus). & D. Placement \\
\rowcolor{nx@mist}
2022 & \citet{kim2022apex} unify sequencer consensus with regional partition tolerance (Apex). & A. Strong \\
2023 & \citet{nakamura2023meridian} propose the first joint placement-and-staleness optimiser evaluated at metro scale (Meridian). & D. Placement \\
\rowcolor{nx@mist}
2024 & Cross-cluster benchmarking proposals appear, but no shared benchmark is adopted (Section~\ref{sec:discussion}). & All \\
\bottomrule
\end{tabularx}
\caption{Representative milestones per year, mapped to the cluster in Figure~\ref{fig:conceptmap}
they belong to.}
\label{tab:timeline}
\end{table}

Two patterns stand out. First, clusters A and B were established almost independently of each other ---
the 2015--2017 papers rarely cite across the strong/eventual boundary, treating the choice as a
binary design decision made once at system-design time. Second, from 2019 onward every major result
either exposes tunability (cluster C) or explicitly couples consistency to placement (cluster D),
suggesting the field's centre of gravity has shifted from \emph{which model to pick} toward \emph{how
to make the model adaptive}.

\begin{keytakeaway}
The field's trajectory runs from picking one consistency model at design time (2015--2018) toward
making consistency an adaptive, per-request, placement-aware decision (2019--2024).
\end{keytakeaway}
